%%%%

%%%   Самарский государственный технический университет

%%%   Кафедра прикладной математики и информатики

%%%

%%%   Преамбула для статьи в журнал "Вестник Самарского государственного

%%%   технического университета. Серия: «Физико-математические науки»"

%%%   Ориентирована под MikTEx 2.4 for Win32

%%%

%%%   Хотя сам журнал собирается под GNU/Linux на дистрибутиве TeXLive



\documentclass[11pt,a4paper,twoside]{article}



\usepackage[cp1251]{inputenc}

\usepackage[T2A]{fontenc}

\usepackage[english,russian]{babel}

\usepackage{samgtubib}

\usepackage[dvips]{graphicx}

\usepackage[multiple]{footmisc}



\usepackage{samgtu}

\renewcommand{\VestnikYear}{2009} % Год издания Вестника

\renewcommand{\VestnikNumber}{2\,(19)} % Номер Вестника

\renewcommand{\VestnikMonth}{Ноябрь} % Месяц выхода Вестника

\renewcommand{\VestnikData}{01.11.09} % Дата подписания Вестника в печать

\renewcommand{\VestnikUPL}{26,64} % Усл. печ. л.

\renewcommand{\VestnikUIL}{25,73} % Уч.-изд. л.

\renewcommand{\VestnikRegNumber}{479/09} % Регистрационный номер

\renewcommand{\VestnikOrder}{994} % Номер заказа









 

\renewcommand{\firstpage}{188}

\renewcommand{\lastpage}{190}





\udk{517.95}

\msc{35P30; 35J60}



\title{О~характере разрывов нелинейности в~задачах на~собственные значения для~уравнений эллиптического типа}{О~характере разрывов нелинейности в~задачах \dots}



\author{Д.\,К.~Потапов}{П\,о\,т\,а\,п\,о\,в~Д.\,К.}



\institute{Санкт-Петербургский государственный университет,\\

198504, Россия, Санкт-Петербург, Университетский просп., 35.}



\email{E-mail}{potapov@apmath.spbu.ru} %E-mail's авторов



\otherinf{\fullauthor{Дмитрий Константинович Потапов} \regalia{к.ф.-м.н., доц.} \position{доцент} \dept{каф. высшей математики}

}



\keywords{краевые задачи, уравнения эллиптического типа, задачи на собственные значения, разрывная нелинейность, характер разрывов}



\workabstract{Рассматриваются задачи на собственные значения для уравнений эллиптического 

типа с разрывными по фазовой переменной нелинейностями. 

Исследуется характер разрывов нелинейности в таких задачах. 

В отличие от работ других авторов, в данной статье ослаблены ограничения на

точки разрыва нелинейности.}



\engtitle{On the character of nonlinearity discontinuities in eigenvalue problems for elliptic equations}





\engauthor{D.\,K.~Potapov}{P\,o\,t\,a\,p\,o\,v~D.\,K.}





\enginstitute{St. Petersburg State University, \\

35, Universitetskiy prosp., St. Petersburg, 198504, Russia.

}



\otherenginf{\fullauthor{Dmitriy K. Potapov} \regalia{Ph. D. (Phys. \& Math.)} \position{Associate Professor}

\dept{Dept. of Higher Mathematics}}



\engabstract{The eigenvalue problems for equations of elliptic type with discontinuous by the phase variable nonlinearities are considered. The character of nonlinearity discontinuities is investigated. In this paper the restrictions on discontinuity points of nonlinearity are weaker than in works of other authors.}



\engkeywords{boundary value problems, elliptic equations, eigenvalue problems, discontinuous nonlinearity, character of discontinuities}







\date{01/IV/2012}

\revisiondate{10/VIII/2012}











\maketitle 



\small



В течение ряда лет автор изучает основные краевые задачи для уравнений

эллиптического типа со спектральным параметром и разрывной по фазовой

переменной нелинейностью (см., например, работы~[\citen{pot1}--\citen{pot32}]).

В работах других отечественных математиков в последние годы нелинейные задачи на 

собственные значения для уравнений эллиптического типа с разрывными 

нелинейностями не рассматривались.      

В работах зарубежных математиков последних лет (см., например,

работы~[\citen{marano1}--\citen{bonanno6}]) на разрывы нелинейности $g(x, u)$ по

фазовой переменной $u$ накладывается следующее достаточно жёсткое ограничение "---

существует множество $\Omega_0\subset\Omega$ меры нуль, для которого 

объединение

$$

\bigcup\limits_{x\in\Omega\backslash\Omega_0}\{u\in\mathbb R: g(x,\cdot)

\mbox{ разрывна в точке } u\}

$$

имеет меру нуль ($\Omega$ "---  область, в которой рассматривается краевая задача).                                                      

По сравнению с работами~[\citen{marano1}--\citen{bonanno6}] в 

работах автора ослаблены ограничения на точки разрыва

нелинейности $g(x,u)$ по $u$, а именно на разрывы функции $g(x,\cdot)$

накладывается условие $g(x,u-)<g(x,u+)$, где $g(x,u\pm)=\lim_{s\to u\pm 0}g(x,s)$.

Таким образом, в работах автора не предполагается, что проекция множества точек разрыва $g(x,u)$ по $u$

на ось фазовой переменной $u$ имеет меру нуль в $\mathbb R$, что достаточно

существенно.

Это важно и для ряда прикладных

задач. Например, в задаче об отрывных течениях несжимаемой жидкости

М.~А.~Гольдштика~\cite{gold} отрывные течения формируются на множестве ненулевой

меры, когда разрывы  прыгающие~\cite{pot5, pot12} (т.~е. если $u$ "---  точка

разрыва функции $g(x,\cdot)$, то $g(x,u-)<g(x,u+)$).

Ограничение на множество точек разрыва нелинейности можно заменить на более

общее "---  А-условие.



\smallskip



{\sc Определение.}

Для дифференциального уравнения с дифференциальным  оператором $L$ {\it выполнено А-условие}, если

найдётся не более чем счётное семейство поверхностей

$\{S_i, i\in I\}$, $S_i=\{(x,u)\in \mathbb R^{n+1} : u=\varphi_i(x), x\in\Omega\}$,

$\varphi_i\in{ W}_{1, \rm loc}^2(\Omega)$,

для которых при почти всех $x\in\Omega$ неравенство $g(x,u-)>g(x,u+)$

влечёт существование $i\in I$ такого, что $u=\varphi_i(x)$ и 

$(L\varphi_i(x)-g(x,\varphi_i(x)-))(L\varphi_i(x)-g(x,\varphi_i(x)+))>0$.



\smallskip



Такое условие предполагалось, например, в работе~\cite{pot7}. 

Отметим, что А-условие запрещает при почти всех $x$ выход решения $u(x)$ краевой задачи на

поверхности разрывов нелинейности $g(x,u)$ в точках так называемых падающих разрывов 

по фазовой переменной, т.~е. точках, для которых $g(x,u-)>g(x,u+)$.

Таким образом, если для почти всех $x\in\Omega$ верно неравенство

$g(x,u-)<g(x,u+)$ $\forall u\in \mathbb R$, т.~е. все разрывы по фазовой 

переменной $u$ "--- прыгающие, то А-условие для уравнения выполняется.

Несложно привести достаточные и легко проверяемые признаки выполнимости

А"=условия, выраженные в терминах коэффициентов дифференциального оператора и

нелинейности, а также дать физическую трактовку этому условию.



\begin{thebibliography}{99}



\RBibitem{pot1}

\by Павленко~В.~Н., Потапов~Д.~К.

\paper О~существовании луча собственных значений для уравнений с~разрывными операторами

\jour Сиб. матем. журн.

\yr 2001

\vol 42

\issue 4

\pages 911--919

\transl англ. пер.:~

\by Pavlenko~V.~N., Potapov~D.~K.

\paper  Existence of a ray of  eigenvalues for equations with discontinuous operators

\jour Siberian Math.~J.

\yr 2001

\vol 42

\issue 4

\pages 766--773





\RBibitem{pot2}

\by Павленко~В.~Н., Потапов~Д.~К.

\paper Аппроксимация краевых задач эллиптического типа со~спектральным параметром и разрывной нелинейностью

\jour Изв. вузов. Матем.

\yr 2005

\issue 4

\pages 49--55

\transl англ. пер.:~

\by Pavlenko~V.~N., Potapov~D.~K. 

\paper Approximation of   boundary-value problems of elliptic type with a spectral parameter and  discontinuous nonlinearity

\jour Russian Math. (Iz. VUZ)

\yr 2005

\vol 49

\issue 4

\pages 46--52



\RBibitem{pot9}

\by Потапов~Д.~К.

\paper Об одной оценке сверху величины бифуркационного  параметра в~задачах на собственные значения для уравнений эллиптического типа  с~разрывными нелинейностями

\jour Дифференц. уравнения

\yr 2008

\vol 44

\issue 5

\pages 715--716

\transl  англ. пер.:~

\by Potapov~D.~K.

\paper On an upper bound for the value of the  bifurcation parameter in eigenvalue problems for elliptic equations with

 discontinuous nonlinearities

 \jour Differ. Equ.

 \yr 2008

 \vol 44

 \issue 5

 \pages 737--739

 

\RBibitem{pot11}

\by Потапов~Д.~К.

\paper О структуре множества собственных значений для уравнений  эллиптического типа высокого порядка с~разрывными нелинейностями

\jour Дифференц. уравнения

\yr 2010

\vol 46

\issue 1

\pages 150--152

\transl  англ. пер.:~

\by Potapov~D.~K.

\paper On the eigenvalue set structure for higher-order   equations of elliptic type with discontinuous nonlinearities

\jour Differ. Equ.

\yr  2010

\vol 46

\issue 1

\pages 155--157



\RBibitem{pot14}

\by Потапов~Д.~К.

\paper Оценки дифференциального оператора в задачах со спектральным  параметром для уравнений эллиптического типа с~разрывными нелинейностями

\jour Вестн. Сам. гос. техн. ун-та. Сер. Физ.-мат. науки

\yr 2010

\issue 5(21)

\pages 268--271

\misctransl

\by Potapov~D.~K.

\paper Estimations of a differential operator in spectral  parameter problems for elliptic equations with discontinuous nonlinearities

\jour  Vestn. Samar. Gos. Tekhn. Univ. Ser. Fiz.-Mat. Nauki

\yr 2010

\issue 5(21)

\pages 268--271



\RBibitem{pot17}

\by Потапов~Д.~К.

\paper Бифуркационные задачи для уравнений эллиптического типа с~разрывными нелинейностями

\jour Матем. заметки

\yr 2011

\vol 90

\issue 2

\pages 280--284

\transl англ. пер.:~

\by Potapov~D.~K.

\paper Bifurcation problems for equations of elliptic type with discontinuous nonlinearities

\jour Math. Notes

\yr 2011

\vol 90

\issue 2

\pages 260--264





\RBibitem{pot32}

\by Потапов~Д.~К.

\paper О~количестве решений в~задачах на собственные значения  для уравнений эллиптического типа с~разрывными нелинейностями

\jour Вестн. Сам. гос. техн. ун-та. Сер. Физ.-мат. науки

\yr 2012

\issue 1(26)

\pages 251--255

\misctransl

\by Potapov~D.~K.

\paper On number of solutions in eigenvalue problems for  elliptic equations with discontinuous nonlinearities

\jour Vestn. Samar.  Gos. Tekhn. Univ. Ser. Fiz.-Mat. Nauki

\yr 2012

\issue 1(26)

\pages 251--255





\Bibitem{marano1}

\by Marano~S.~A.

\paper Elliptic eigenvalue problems with highly discontinuous  nonlinearities

\jour Proc. Amer. Math. Soc.

\yr 1997

\vol 125

\issue 10

\pages 2953--2961



\Bibitem{marano2}

\by Marano~S.~A., Motreanu~D.

\paper On a three critical points theorem for  non-differentiable functions and applications to nonlinear boundary value  problems

\jour Nonlinear Anal.

\yr 2002

\vol 48

\issue 1

\pages 37--52



\Bibitem{bonanno2}

\by Bonanno~G.

\paper Some remarks on a three critical points theorem

\jour Nonlinear Anal.

\yr 2003

\vol 54

\issue 4

\pages 651--665



\Bibitem{bonanno3}

\by Bonanno~G., Giovannelli~N.

\paper An eigenvalue Dirichlet problem involving  the $p$-Laplacian with discontinuous nonlinearities

\jour J.~Math. Anal. Appl.

\yr 2005

\vol 308

\issue 2

\pages 596--604



\Bibitem{zhang}

\by Zhang~G., Liu~S.

\paper Three symmetric solutions for a class of  elliptic involving the $p$-Laplacian with discontinuous nonlinearities in  $\mathbb R^n$

\jour Nonlinear Anal.

\yr 2007

\vol 67

\issue 7

\pages 2232--2239



\Bibitem{bonanno4}

\by  Bonanno~G., Candito~P.

\paper Non-differentiable functionals and applications  to elliptic problems with discontinuous nonlinearities

\jour J.~Differential Eq.

\yr 2008

\vol 244

\issue 12

\pages 3031--3059



\Bibitem{bonanno6}

\by Bonanno~G., Chinni~A.

\paper Discontinuous elliptic problems involving the   $p(x)$-Laplacian

\jour Math. Nachr.

\yr 2011

\vol 284

\issue 5--6

\pages 639--652



\RBibitem{gold}

\by Гольдштик~М.~А.

\paper Математическая модель отрывных течений несжимаемой  жидкости

\jour Докл. АН СССР

\yr 1962

\vol 147

\issue 6

\pages 1310--1313

\transl  англ. пер.:~

\by Gol'dshtik~M.~A.

\paper A mathematical model of separated flows  in an incompressible liquid

\jour Soviet. Math. Dokl.

\yr 1963

\vol 7

\pages 1090--1093





\RBibitem{pot5}

\by Потапов~Д.~К.

\paper Математическая модель отрывных течений несжимаемой  жидкости

\jour Изв. РАЕН. Сер.~МММИУ

\yr 2004

\vol 8

\issue 3--4

\pages 163--170

\misctransl 

\by Potapov~D.~K.

\paper A mathematical model of separated flows in an  incompressible fluid

\jour Izv. RAEN. Ser.~MMMIU

\yr 2004

\vol 8

\issue 3--4

\pages 163--170



\RBibitem{pot12}

\by Потапов~Д.~К.

\paper Непрерывные аппроксимации задачи Гольдштика

\jour Матем. заметки

\yr 2010

\vol 87

\issue 2

\pages 262--266

\transl  англ. пер.:~

\by Potapov~D.~K.

\paper Continuous approximations of Gol'dshtik's model

\jour  Math. Notes

\yr 2010

\vol 87

\issue 2

\pages 244--247



\RBibitem{pot7}

\by Потапов~Д.~К.

\paper Аппроксимация задачи Дирихле для уравнения  эллиптического типа высокого порядка со спектральным параметром и~ разрывной нелинейностью

\jour Дифференц. уравнения

\yr 2007

\vol 43

\issue 7

\pages 1002--1003

\transl  англ. пер.:~

\by Potapov~D.~K.

\paper Approximation to the Dirichlet problem for a  higher-order elliptic equations with a spectral parameter and a discontinuous nonlinearity

 \jour Differ. Equ.

 \yr 2007

 \vol 43

 \issue 7

 \pages 1031--1032



\end{thebibliography}



\noindent

\hskip6.5cm \thedate \par\noindent

\hskip6.5cm\therevdate





\vspace{-5mm}





\makeengtitle



\newpage



%\end{document}

